\documentclass[12pt,a4paper]{article}

\usepackage[spanish,es-nodecimaldot]{babel}
\usepackage[utf8]{inputenc}
\usepackage[T1]{fontenc}
\usepackage{lmodern}
\usepackage{geometry}
\usepackage{microtype}
\usepackage{setspace}
\usepackage{parskip}
\usepackage{enumitem}
\usepackage{xcolor}
\usepackage{listings}
\usepackage{hyperref}
\usepackage{amsmath}
\usepackage{amssymb}
\usepackage{array}
\usepackage{booktabs}
\usepackage{longtable}
\usepackage{tcolorbox}

\geometry{margin=2.5cm}
\onehalfspacing
\setlist[itemize]{leftmargin=*,itemsep=0.2em}
\setlist[enumerate]{leftmargin=*,itemsep=0.2em}

\hypersetup{
  colorlinks=true,
  linkcolor=black,
  urlcolor=blue,
  pdftitle={Resumen teorico TP5 - Cobertura logica},
  pdfauthor={}
}

\newtcolorbox{definicionbox}[1][]{
  colback=blue!4!white,
  colframe=blue!50!black,
  fonttitle=\bfseries,
  title=Definicion,
  #1
}

\newtcolorbox{ideabox}[1][]{
  colback=green!4!white,
  colframe=green!40!black,
  fonttitle=\bfseries,
  title=Idea clave,
  #1
}

\newtcolorbox{ejemplobox}[1][]{
  colback=orange!4!white,
  colframe=orange!50!black,
  fonttitle=\bfseries,
  title=Ejemplo,
  #1
}

\lstdefinestyle{codestyle}{
  basicstyle=\ttfamily\small,
  frame=single,
  breaklines=true,
  showstringspaces=false
}

\title{\textbf{Resumen Teorico -- Practico 5}\\
Testing basado en expresiones logicas\\
Predicados, clausulas y cobertura logica\\
\textit{Testing de Software 2024}}
\author{}
\date{}

\begin{document}

\maketitle
\tableofcontents
\newpage

\section*{Introduccion}
\addcontentsline{toc}{section}{Introduccion}

Este resumen sintetiza el material teorico asociado al TP5:

\begin{itemize}
  \item \textit{Introduction to Software Testing} (Ammann \& Offutt, 2da edicion), Capitulo 8:
  \textbf{Logic Coverage}.
  \item \textit{Notas 09: Cobertura logica} (V. Bengolea, primera parte).
  \item \textit{Notas 09: Cobertura logica} (V. Bengolea, segunda parte).
\end{itemize}

\textbf{Idea general.} El software contiene constantemente \emph{expresiones logicas}: las
guardas de un \texttt{if} o \texttt{while}, las transiciones de una FSM, las condiciones
mencionadas en una especificacion o un requisito. Una expresion logica con $n$ variables
booleanas tiene $2^n$ combinaciones; en la practica es inviable probarlas todas. Los
\textbf{criterios de cobertura logica} eligen un subconjunto pequeno y representativo, atando
los tests a la estructura de la formula.

\begin{ideabox}
Los criterios de cobertura logica responden a una pregunta operativa: ``si una decision booleana
toma muchos valores distintos, \textquestiondown{}cuales necesito probar para estar razonablemente seguro de que
la condicion fue codificada bien?''.
\end{ideabox}

\section{Predicados y clausulas}

\subsection{Definiciones basicas}

\begin{definicionbox}
\textbf{Predicado.} Expresion que se evalua a un valor booleano. Un predicado puede contener:
\begin{itemize}
  \item variables booleanas,
  \item variables no booleanas combinadas con operadores relacionales ($<, ==, !=, \ldots$),
  \item llamadas a funciones que devuelven un booleano.
\end{itemize}
La estructura interna se forma con los operadores logicos
$\neg,\;\wedge,\;\vee,\;\rightarrow,\;\oplus,\;\leftrightarrow$.

\textbf{Clausula.} Predicado que \emph{no contiene operadores logicos}. En la formula
$((a > b) \vee C) \wedge p(x)$ las clausulas son $(a>b)$, $C$ y $p(x)$.
\end{definicionbox}

\begin{ideabox}
Las \emph{clausulas} son los ``atomos'' booleanos de una formula: son lo mas pequeno que aun se
evalua a \textit{true} o \textit{false}. Los criterios actuan sobre clausulas porque cada
clausula corresponde, conceptualmente, a una decision atomica que el programador codifica.
\end{ideabox}

\subsection{\textquestiondown{}De donde salen los predicados?}

Las expresiones logicas no aparecen solo en el codigo. Surgen tambien de:

\begin{itemize}
  \item \textbf{Codigo:} guardas de \texttt{if}/\texttt{while}, condiciones de salida de bucles,
  cortocircuitos.
  \item \textbf{Maquinas de estados y statecharts:} las transiciones se etiquetan con guardas
  booleanas.
  \item \textbf{Requisitos y especificaciones} (formales o informales). Por ejemplo, las
  precondiciones de un metodo: ``the stack is not full and o is not null'' se traduce a
  $\neg\,\textit{this.isFull}() \wedge o \ne \textit{null}$.
  \item \textbf{Lenguaje natural,} convirtiendo enunciados a logica formal.
  \item \textbf{Consultas SQL} y otras DSL que evaluan condiciones.
\end{itemize}

La FAA estadounidense exige cobertura de expresiones logicas (concretamente MCDC, un primo
hermano de CACC) en \textit{safety-critical software}.

\subsection{Traduccion de lenguaje natural}

Pasar de una frase a una formula logica es una fuente clasica de errores:

\begin{itemize}
  \item \emph{``I am interested in SWE 637 and CS 652''}: la conjuncion natural es realmente un
  \emph{O} logico: $\textit{course} = \textit{swe637} \vee \textit{course} = \textit{cs652}$.
  \item \emph{``If you leave before 6:30, take Braddock; if you leave after 7:00, take
  Prosperity''}: la formula directa $(t < 6{:}30 \rightarrow \textit{Braddock}) \wedge
  (t > 7{:}00 \rightarrow \textit{Prosperity})$ es incompleta: deja sin cubrir el intervalo
  $[6{:}30, 7{:}00]$. Una version corregida es $(t < 6{:}30 \rightarrow \textit{Braddock})
  \wedge (t \ge 6{:}30 \rightarrow \textit{Prosperity})$.
\end{itemize}

\begin{ideabox}
\textbf{``Y'' del castellano} a menudo se traduce a \emph{O logico}; \textbf{``O'' del
castellano} suele significar XOR (``cafe o te''); las descripciones por casos suelen tener
\emph{huecos} no cubiertos. Antes de aplicar criterios de cobertura conviene asegurarse de que
la traduccion sea fiel.
\end{ideabox}

\subsection{Datos empiricos}

Estudio sobre 63 programas open source y mas de $400\,000$ predicados:

\begin{itemize}
  \item $88{,}5\%$ tiene una sola clausula.
  \item $9{,}5\%$ tiene dos.
  \item $1{,}35\%$ tiene tres.
  \item $0{,}65\%$ tiene cuatro o mas.
\end{itemize}

Es decir, los predicados con muchas clausulas son la minoria, pero son justamente donde la
diferencia entre criterios (PC, CC, CACC, \ldots) se vuelve interesante.

\subsection{Expresiones logicas desde el codigo}

Una decision como
\begin{lstlisting}[style=codestyle,language=Java]
if (((a > b) || C) && (x < y)) o.m(); else o.n();
\end{lstlisting}
produce el predicado $p = ((a>b) \vee C) \wedge (x<y)$ con clausulas $(a>b)$, $C$ y $(x<y)$.

\section{Notacion estandar}

\begin{definicionbox}
A lo largo del capitulo se usa la siguiente notacion:
\begin{itemize}
  \item $P$: conjunto de predicados.
  \item $p$: un predicado de $P$.
  \item $C$: conjunto de clausulas en $P$.
  \item $C_p$: conjunto de clausulas del predicado $p$.
  \item $c$: una clausula.
  \item $\textit{RT}$: conjunto de requisitos de test generados por el criterio.
\end{itemize}
\end{definicionbox}

\section{Cobertura de predicados y de clausulas}

Los dos criterios mas simples piden que cada predicado o cada clausula tome ambos valores de
verdad.

\begin{definicionbox}
\textbf{Predicate Coverage (PC).} Para cada predicado $p \in P$, $RT$ contiene dos requisitos:
$p$ se evalua a \textit{true} y $p$ se evalua a \textit{false}.
\end{definicionbox}

\begin{definicionbox}
\textbf{Clause Coverage (CC).} Para cada clausula $c \in C$, $RT$ contiene dos requisitos: $c$
se evalua a \textit{true} y $c$ se evalua a \textit{false}.
\end{definicionbox}

\textbf{Relacion con cobertura estructural.} Si los predicados provienen de las decisiones del
codigo, entonces:

\begin{itemize}
  \item PC equivale a \emph{cobertura de aristas} (Edge Coverage) del CFG: cada rama del
  \texttt{if} tiene que ejecutarse.
  \item CC equivale a \emph{cobertura de condiciones}: cada condicion booleana individual debe
  evaluar a $T$ y a $F$.
\end{itemize}

\begin{ideabox}
\textbf{CC y PC no se subsumen entre si.} PC ignora la estructura interna del predicado; CC
ignora si el predicado toma realmente ambos valores. Pueden construirse suites que cumplan uno
sin cumplir el otro.
\end{ideabox}

\subsection{Ejemplo: \texorpdfstring{$p = a \vee (b \wedge c)$}{p = a or (b and c)}}

\begin{ejemplobox}
Suite $\{(T,F,F), (F,T,T)\}$:
\begin{itemize}
  \item Cada clausula toma $T$ y $F$ $\Rightarrow$ CC ok.
  \item Pero $p$ vale $T$ en los dos casos $\Rightarrow$ PC \emph{no} se cumple.
\end{itemize}

Suite $\{(T,T,T), (F,T,F)\}$:
\begin{itemize}
  \item $p$ toma $T$ y $F$ $\Rightarrow$ PC ok.
  \item La clausula $b$ queda fija en $T$ $\Rightarrow$ CC \emph{no} se cumple.
\end{itemize}
\end{ejemplobox}

\subsection{Limitacion practica de CC}

CC pide que cada clausula tome ambos valores, pero \emph{sin coordinarse} con el predicado. En
$p = a \wedge b$, los pares $(a=T,b=F)$ y $(a=F,b=T)$ satisfacen CC, pero ninguno ejercita el
caso $p=T$. Por eso CC suele ir acompanado de PC o reemplazado por criterios mas finos como
CACC.

\section{Cobertura combinatoria}

\begin{definicionbox}
\textbf{Combinatorial Coverage (CoC).} Para cada predicado $p \in P$, $RT$ contiene un requisito
por cada combinacion posible de valores de las clausulas de $C_p$.
\end{definicionbox}

CoC es ``el criterio ideal'': cubrir toda la tabla de verdad. El costo es exponencial: con $n$
clausulas hace falta $2^n$ tests. En la practica se reserva solo para predicados muy chicos.

\textbf{Cuello de botella.} CoC subsume PC y CC, pero suele ser inviable para
$n \ge 4$ o cuando hay muchos predicados.

\section{Clausulas activas}

CC peca por desacoplado y CoC por excesivo. La idea intermedia es: para cada clausula, exigir
casos en los que esa clausula realmente \emph{determine} el valor del predicado.

\subsection{Determinacion}

\begin{definicionbox}
Una clausula $c_i \in C_p$ \textbf{determina} el predicado $p$ si, al fijar las demas clausulas
(las menores) y cambiar $c_i$, el valor de $p$ cambia. De forma equivalente, $p$ se vuelve una
funcion de $c_i$ solo: $p|_{c_i=T} \neq p|_{c_i=F}$.
\end{definicionbox}

\begin{ejemplobox}
En $p = a \vee b$, la clausula $a$ determina cuando $b=F$ (sino el OR siempre es $T$). En
$p = a \wedge b$, $a$ determina cuando $b=T$ (sino el AND siempre es $F$).
\end{ejemplobox}

\subsection{Active Clause Coverage (ACC)}

\begin{definicionbox}
\textbf{Active Clause Coverage (ACC).} Para cada $p \in P$ y cada clausula $c_i \in C_p$ (la
\emph{clausula mayor}), elegir valores para las clausulas menores $c_j$ ($j \ne i$) tales que
$c_i$ determina $p$. $RT$ contiene dos requisitos por clausula mayor:
\begin{enumerate}
  \item $c_i = T$ con las menores fijadas asi;
  \item $c_i = F$ con las menores fijadas asi.
\end{enumerate}
\end{definicionbox}

ACC se queda ambiguo porque no precisa \emph{como} se eligen las clausulas menores. De ahi
nacen tres variantes que se diferencian en la rigidez de esa eleccion.

\subsection{GACC, CACC y RACC}

\begin{definicionbox}
\textbf{General Active Clause Coverage (GACC).} Para cada $c_i$, las menores pueden tener
cualquier valor mientras $c_i$ determine $p$. No se exige nada sobre el valor de $p$ entre los
dos tests.

\textbf{Correlated Active Clause Coverage (CACC).} Igual que GACC, pero \emph{ademas} se exige
que en el par de tests $p$ tome ambos valores ($T$ y $F$): la clausula mayor debe \emph{causar
correlativamente} el cambio.

\textbf{Restricted Active Clause Coverage (RACC).} Igual que CACC, pero las clausulas menores
deben tener \emph{exactamente los mismos valores} en los dos tests del par.
\end{definicionbox}

\begin{ideabox}
\textbf{De debil a fuerte:} $\text{GACC} \prec \text{CACC} \prec \text{RACC}$. CACC es el mas
usado en la practica; es esencialmente lo que la industria llama \textbf{MCDC} (Modified
Condition / Decision Coverage). RACC es mas exigente y a veces \emph{infactible} (no existen
asignaciones validas que dejen las menores fijas y permitan que la mayor determine).
\end{ideabox}

\subsection{Ejemplo: \texorpdfstring{$p = a \wedge (b \vee c)$}{p = a and (b or c)}}

\begin{ejemplobox}
\textbf{Mayor $a$.} $a$ determina cuando $b \vee c = T$. Par CACC:
$(T,T,F) \Rightarrow p=T$ vs $(F,T,F) \Rightarrow p=F$.

\textbf{Mayor $b$.} $b$ determina cuando $c=F$ y $a=T$. Par CACC:
$(T,T,F) \Rightarrow p=T$ vs $(T,F,F) \Rightarrow p=F$.

\textbf{Mayor $c$.} $c$ determina cuando $b=F$ y $a=T$. Par CACC:
$(T,F,T) \Rightarrow p=T$ vs $(T,F,F) \Rightarrow p=F$.

Las menores estan fijadas a un unico valor por par, asi que esta suite es RACC tambien.
\end{ejemplobox}

\section{Cobertura inactiva}

La idea simetrica: para cada clausula mayor, mantenerla en cada valor mientras \emph{no}
determina el predicado, y mostrar que en ese caso $p$ no depende de $c_i$.

\begin{definicionbox}
\textbf{Inactive Clause Coverage (ICC).} Para cada $c_i$ se eligen menores que hagan que $c_i$
\emph{no} determine $p$, y se exigen cuatro requisitos por clausula mayor:
\begin{enumerate}
  \item $c_i = T$, $p = T$;
  \item $c_i = T$, $p = F$;
  \item $c_i = F$, $p = T$;
  \item $c_i = F$, $p = F$.
\end{enumerate}
\textbf{GICC} no restringe las menores; \textbf{RICC} las pide identicas en cada par.
\end{definicionbox}

ICC ayuda a detectar errores de tipo ``clausula sobrante'': si el codigo agrega una clausula
$c_i$ que en la formula correcta no aparece, ICC trata de exhibirlo verificando que $p$ se
mantiene igual cuando $c_i$ cambia ``por las razones correctas''.

\section{Mapa de subsuncion}

Las relaciones de subsuncion entre los criterios (un criterio $X$ \emph{subsume} a $Y$ si toda
suite que satisface $X$ tambien satisface $Y$) son:

\begin{itemize}
  \item $\text{CoC} \succ \text{CACC} \succ \text{GACC} \succ \text{CC}$.
  \item $\text{CACC} \succ \text{PC}$.
  \item $\text{RACC} \succ \text{CACC}$, pero $\text{RACC}$ puede ser \emph{infactible} en
  ciertos predicados.
  \item Ni CC subsume a PC, ni PC subsume a CC.
\end{itemize}

\begin{ideabox}
Una manera intuitiva de pensar la jerarquia:
\begin{itemize}
  \item PC ve el predicado como una caja negra: ``salgan ambos resultados''.
  \item CC mira las clausulas, pero sin coordinarlas con $p$.
  \item ACC/CACC/RACC son lo que se quiere casi siempre: cada clausula tiene que importar.
  \item CoC es la cota superior brutal: probar todo.
\end{itemize}
\end{ideabox}

\section{Origenes de las expresiones logicas}

\subsection{Desde el codigo fuente}

Una expresion booleana del codigo (sin importar como esta agrupada con parentesis) se mapea a
una formula \emph{con clausulas} en la que cada operando relacional o booleano elemental es una
clausula. El cortocircuito (\texttt{\&\&}, \texttt{||}) en lenguajes como Java suele tratarse
como evaluacion sintactica completa para fines del criterio.

\subsection{Desde especificaciones}

Las precondiciones y postcondiciones (formato JML, contratos, comentarios estilo Eiffel)
contienen formulas logicas explicitas. Por ejemplo:

\begin{lstlisting}[style=codestyle]
/*@precondition: the stack is not full and o object is not null; */
public void push(Object o);
\end{lstlisting}
se traduce a $\neg\,\textit{this.isFull}() \wedge o \ne \textit{null}$, con dos clausulas claras
sobre las que aplicar PC/CC/CACC.

\subsection{Desde maquinas de estados}

Las transiciones de un statechart estan etiquetadas por \emph{guardas} de la forma
$\textit{trigger}\,[\textit{guard}] / \textit{accion}$. La guarda es un predicado al que se le
puede pedir cobertura logica: PC sobre transiciones, CACC sobre guardas con varias clausulas.

\subsection{Desde lenguaje natural}

Como se discutio, la traduccion es propensa a errores de cuantificadores, conectivas y casos no
cubiertos. Antes de aplicar cobertura logica conviene formalizar y revisar la formula.

\section{Conexion con MCDC}

Cuando la industria (especialmente aviacion, DO-178B/C) habla de cobertura logica usualmente se
refiere a \textbf{MCDC}: \emph{Modified Condition/Decision Coverage}. Es esencialmente lo mismo
que \textbf{CACC}: cada decision toma todos los valores; cada condicion toma todos los valores;
y se demuestra que cada condicion influye de forma independiente en el resultado de la decision.

\begin{ideabox}
Si te piden ``MCDC'' en una norma como DO-178C, la receta operativa es la misma que para CACC:
para cada clausula mayor, dos casos donde ella determina el predicado y donde el predicado
cambia.
\end{ideabox}

\section{Como derivar tests con cobertura logica}

\subsection{Receta paso a paso}

\begin{enumerate}
  \item Identificar el predicado $p$ y sus clausulas $C_p$.
  \item Construir o pensar la tabla de verdad (o un subconjunto razonado de ella).
  \item Elegir el criterio (PC, CC, CACC, RACC, \ldots) segun el nivel de exigencia.
  \item Para cada requisito del criterio, encontrar valores de las clausulas que lo satisfagan.
  \item Traducir cada combinacion de clausulas a inputs concretos del codigo (a veces hace falta
  invertir las relaciones $a = (x < y)$ para encontrar $x,y$).
  \item Escribir los tests y verificar con la herramienta correspondiente.
\end{enumerate}

\subsection{De clausulas a inputs}

Una clausula como $a := \textit{curTemp} < \textit{dTemp} - \textit{thresholdDiff}$ es
verdadera cuando se cumple la desigualdad. Para forzar $a = T$ se eligen valores numericos que
la cumplan. Cuando varias clausulas comparten variables (e.g. \texttt{curTemp} aparece en $a$ y
$c$), hay que armar inputs que respeten todas las restricciones simultaneamente; a veces es
necesario fijar parametros adicionales (\texttt{thresholdDiff}, \texttt{minLag}) para tener
control fino.

\section{Pequeno glosario}

\begin{itemize}
  \item \textbf{Clausula mayor.} La clausula sobre la que se esta razonando (la que ``cambia'')
  en un par CACC/RACC.
  \item \textbf{Clausulas menores.} Las restantes; se eligen para que la mayor determine $p$.
  \item \textbf{Test factible.} Aquel para el que existe una entrada del programa que produzca
  esa combinacion de valores. Si la combinacion es contradictoria con la logica del codigo
  (e.g.\ dos clausulas se excluyen), el test es \emph{infactible}.
  \item \textbf{Subsuncion.} Si $X$ subsume a $Y$, todo conjunto de tests que satisface $X$
  satisface tambien $Y$.
\end{itemize}

\section*{Conclusion}
\addcontentsline{toc}{section}{Conclusion}

La cobertura logica organiza, con definiciones claras, una intuicion vieja: ``si una decision
del programa tiene varias condiciones, hay que mover cada una y mirar como cambia el resultado''.
PC y CC son puntos de partida, CACC/RACC formalizan la idea de \emph{causalidad}, y CoC actua
como cota superior. En la practica, la combinacion mas usada es CACC (equivalente a MCDC), por
ser exigente sin volverse exponencial.

\end{document}
